
Now we recall some basic definitions about graphs and simplicial complexes that will be useful.

Let $G$ be a graph with vertex set $V(G)$
and edge set $E(G)$.  
A subgraph $H \subseteq G$  is called \textit{induced} if $\{u,v\}$ is an edge of $H$ if and only if $u$ and $v$ are vertices of $H$ and $\{u,v\}$ is an edge of $G$. 
A \textit{clique} in a graph is an induced subgraph that is a  \emph{complete graph}.

For $u\in V(G)$, let $N_G(u) = \{v \in V (G)\mid \{u, v\} \in E(G)\}$ and $N_G[u]=N_G(u)\cup \{u\}$.
For $U \subseteq V(G)$, denote by $G \setminus U$
the induced subgraph of $G$ on the vertex set $V(G) \setminus U$.

Let $G$ be a graph. We denote the graph consisting of two disjoint edges
 by $2K_2$.
%  in $G$ if $G$ does not have an edge with one
%endpoint in $f_1$ and the other in $f_2$.
A graph without induced copy of $2K_2$ is called $2K_2$-free or \textit{gap-free} graph. The \textit{complement} of a graph $G$, denoted by $G^c$, is the graph on the same
vertex set in which $\{u,v\}$ is an edge of $G^c$ if and only if it is not an edge of $G$.
Then  $G$ is gap-free if and only if $G^c$ contains no induced $4$-cycle.
%Thus, $G$ is gap-free if and only if it does not contain two vertex-disjoint edges as an induced subgraph.

%\par A \textit{matching} in a graph $G$ is a subgraph consisting of pairwise disjoint edges.
%If the
%subgraph is an induced subgraph, the matching is an \textit{induced matching}.
%T%he largest size of an induced matching in $G$ is called its induced
%matching number and denoted by $\nu(G)$.
%Let $H$ be a subgraph of $G$. The largest size of induced matching of $H$ as well as
%$G$ denoted by $\nu_{GH}$.  A \textit{dominating induced matching} of $G$ is an induced matching which also
%forms a maximal matching of $G$
A graph $G$ is chordal (also called triangulated) if every induced
cycle in $G$ has length $3$, and is co-chordal if the complement graph $G^c$ is chordal.
 The following important theorem(s) characterizes the edge ideals with regularity 2, and the regularity of their powers.

 

\begin{theorem}
\begin{enumerate}
\item \textup{(}Fr\"oberg~\cite[Thm.1]{froberg}\textup{)} For any graph $G$, we have $G^c$ is chordal if and only if $\reg(I(G))=2$. 

\item \textup{(}Herzog--Hibi--Zheng~\cite[Thm.1.2]{HHZ}\textup{)} Further, in this case $\reg(I(G)^s)=2s$ for any $s$.
\end{enumerate}
\label{thm:chordal}
\end{theorem}

  A \emph{simplicial complex} $\Delta$ on a vertex set $\{1,\ldots, n\}$ is a collection of subsets of $\{1,\ldots ,n\}$ such that if $\tau \in \Delta, \sigma \subseteq \tau$ then we have $\sigma \in \Delta$ and all singletons belong to that collection; if $\tau$ is such a subset belonging to $\Delta$ then $\tau$ is called a \emph{face} of $\Delta$. 
 The \emph{induced subcomplex} $\Delta [A]$ of $\Delta$ on vertex set $A\subset \{1,\ldots,n\}$ is the collection of faces $\tau'$ of $\Delta$ such that $\tau' \subset A$. 
 Clearly the induced subcomplex is a simplicial complex itself. We denote by $\Vr(\Delta)$ the set of vertices of $\Delta$.

 For sets  $A$ and $B$ we sometimes write $A\cup_C B$ for the set $A\cup B$ as a shorthand for denoting their intersection by $C=A\cap B$. 
 Likewise for simplicial complexes.

 The \emph{link} of a vertex $v$ in $ \Delta$ is $$\text{link}_v (\Delta)= \{ \tau |\tau \cup \{v\} \in \Delta, \{v\} \cap \tau =\emptyset\}.$$

 The (open) \emph{star} of a vertex $v$ in a simplicial complex $\Delta$ is the set of all faces that contain $v$, namely $\St_v(\Delta)=\{\tau| \tau \in \Delta, v \in \tau\}$. The \emph{closed star} $\bar{\St}_v (\Delta$ ) of $v$ is defined by the smallest subcomplex that contains $\St_v(\Delta)$. The \emph{antistar} of vertex $v$ is defined as the subcomplex $\text{ast}_v (\Delta) = \{\tau \in \Delta| \tau \cap \{v\} = \emptyset\}$.


 The \emph{join} of two simplicial complexes $\Delta_1$ and $\Delta_2$ is defined by $\Delta_1 * \Delta_2=\{\sigma \cup \tau| \sigma \in \Delta_1, \tau \in \Delta_2\}$. 
 The \emph{suspension} of a simplicial complex $\Delta$, w.r.t. two points $a$ and $b$ not vertices of $\Delta$, is the join defined by $\Sigma_{a,b} \Delta=\Delta * \{\{a\},\{b\},\emptyset \}$; its geometric realization is homeomorphic to the topological suspension of the space $\Delta$.

  For $H$ a graph
  %and $\Delta=\cl(G^c)$, where
  let $\cl(H)$ denote its clique complex, i.e. the simplicial complex whose faces are the cliques of $H$.

The following formulation of regularity follows from the so called Hochster's formula
%and Stanley-Reisner ideal formulation
(see {\cite{ms2005}} for further details):

\begin{theorem}[Hochster's formula]
For every graph $G$ whose edge set is nonempty and $\Delta=\cl(G^c)$ we have:
$$\reg(G)\coloneqq\reg(I(G))= \max\{l+2:\exists W \subseteq \Vr(\Delta), \widetilde{H}_l(\Delta[W];k)\neq 0\}.$$
\end{theorem}
(Sometimes Hochster's formula is phrased with the extra condition on $W$ that $|W|=l+2+i$ for some $i\ge 0$. However, this condition is superfluous: if $|W|\le l+1$ then either $\dim(\Delta[W])<l$ or $\Delta[W]$ is the simplex on $l+1$ vertices. In both cases  $\widetilde{H}_l(\Delta[W])=0$.)
